% ============================================================================
% Section 10: Discussion
% ============================================================================

\section{Discussion}
\label{sec:discussion}

\subsection{Threats to Validity}
\label{sec:discussion:threats}

\paragraph{Internal Validity.}
Our stochastic test semantics assume that trial outcomes $r_1, \ldots,
r_n$ are independent and identically distributed (i.i.d.) draws from a
Bernoulli distribution with parameter $p$. This assumption may be
violated in practice due to:
\begin{enumerate}[leftmargin=*]
  \item \emph{Model caching}: some LLM providers cache responses for
    identical inputs, reducing variance artificially.
  \item \emph{Rate limiting}: sequential API calls may encounter
    different latency profiles, affecting tool-call-dependent agents.
  \item \emph{Temporal drift}: the underlying model may be updated during
    a test run, violating stationarity.
\end{enumerate}
We mitigate (1) by varying temperature and seed parameters across
trials, (2) by randomizing trial order with jitter, and (3) by
recommending test windows shorter than typical model update cycles.
Future work should formalize robustness to non-i.i.d.\ settings.

\paragraph{External Validity.}
Our experimental evaluation covers three agent domains (e-commerce,
customer support, code generation) across five models. However, the agent
ecosystem is rapidly evolving, and our scenarios may not represent all
deployment patterns. In particular, we do not evaluate:
\begin{itemize}[leftmargin=*]
  \item Long-horizon agents with hundreds of steps
  \item Agents with physical-world actuators (robotics, IoT)
  \item Adversarial settings where agents face deliberate manipulation
\end{itemize}
These represent important directions for future generalization.

\paragraph{Construct Validity.}
The mutation adequacy theorem (\cref{thm:mutation-adequacy}) relies on the
coupling effect assumption---that complex faults are coupled to simple
mutations. While the coupling effect is well-established for source code
mutations~\citep{jia2011mutation}, its applicability to agent mutations
(prompt perturbations, tool modifications) has not been empirically
validated at scale. Our experiments (\cref{sec:experiments:characterization}) provide
initial evidence but do not constitute a definitive validation.

\subsection{Limitations}
\label{sec:discussion:limitations}

\paragraph{Cost of Stochastic Testing.}
Running $n$ trials per scenario incurs $n\times$ the cost of a single
execution. For expensive models (e.g., GPT-4o at \$5--15 per complex
agent run), a test suite of 50 scenarios at $n=30$ trials requires 1,500
agent invocations. The token-efficient testing framework
(\cref{sec:token-efficient}) addresses this with 5--20$\times$ cost
reductions through behavioral fingerprinting, adaptive budgets,
trace-first analysis, multi-fidelity proxy testing, and warm-start
sequential analysis. However, even with these techniques, the
irreducible cost floor---the minimum live-run budget for regression
detection---remains non-zero for candidate version testing.
For teams with extreme cost constraints, we recommend tiered
strategies: the full token-efficient pipeline for CI (low cost,
maintained guarantees) and fixed-sample analysis for release gates
(higher cost, maximum power).

\paragraph{Evaluator Reliability.}
When the evaluator $E$ is model-based (e.g., an LLM judge), it
introduces a second source of stochasticity. The true pass probability
becomes $\Pr[\text{agent correct}] \times \Pr[\text{evaluator correct}]$,
which conflates agent quality with evaluator quality. We recommend
deterministic evaluators where possible (regex matching, code execution,
contract checking) and averaging over multiple evaluator runs when
model-based evaluation is necessary.

\paragraph{State-Space Coverage Calibration.}
The $\lambda$ parameter in state-space coverage
(\cref{def:state-coverage}) requires calibration to the agent class.
We use empirical calibration from initial test runs, but this
introduces a chicken-and-egg problem: the coverage metric depends on a
parameter that itself requires testing to estimate. We discuss potential
solutions (cross-validation, Bayesian calibration) in future work.

\paragraph{Equivalent Mutant Detection.}
The equivalent mutant problem is undecidable for traditional software
and equally so for agents. Our heuristic for presuming equivalence
(no significant difference after $n_{\text{equiv}}$ trials on
$k_{\text{equiv}}$ scenarios) may misclassify non-equivalent mutants with
subtle behavioral differences, inflating the mutation score. More
sophisticated equivalence detection (e.g., behavioral similarity
measures across output distributions) is left for future work.

\subsection{Implications of Token-Efficient Testing}
\label{sec:discussion:token-implications}

\paragraph{Enterprise Adoption.}
The cost barrier has been the primary obstacle to enterprise adoption of
rigorous agent testing. Our experiments (\cref{sec:experiments:e7})
demonstrate that the full token-efficient pipeline reduces the cost of a
regression check from hundreds of dollars to single digits, bringing
statistical agent testing within the budget of standard CI/CD
pipelines. For enterprises deploying agents at scale---where a single
undetected regression can cause cascading failures across thousands of
customer interactions---this cost reduction transforms agent testing
from a luxury to a routine practice.

\paragraph{Comparison with Monitoring-Only Approaches.}
Current industry practice relies heavily on production monitoring (e.g.,
LangSmith traces, custom dashboards) to detect agent regressions
\emph{after deployment}. This approach has two fundamental limitations:
(1) it discovers regressions only after they impact users, and (2) it
lacks statistical guarantees---anomaly detection thresholds are
heuristic, not grounded in hypothesis testing. \agentassay{}'s
token-efficient testing enables \emph{pre-deployment} regression
detection with formal guarantees, complementing (not replacing)
production monitoring. The trace-first analysis pillar
(\cref{sec:token:trace-first}) bridges the gap: production traces
collected by monitoring systems can be directly consumed by
\agentassay{} for offline coverage analysis, contract checking, and
metamorphic testing at zero additional cost.

\paragraph{Behavioral Fingerprints as a Representation.}
The behavioral fingerprint (\cref{def:fingerprint}) has applications
beyond regression testing. Fingerprints can serve as agent
\emph{identity signatures} for detecting model drift in production,
as features for automated agent classification (e.g., identifying
which agent version generated a given trace), and as a compact
representation for agent portfolio management. The low effective
dimension $d_{\mathrm{eff}}$ observed in our experiments confirms that
agent behavior admits a compressed representation, with implications
for agent monitoring, debugging, and optimization.

\paragraph{Continuous Assay in Production.}
The combination of trace-first analysis and warm-start SPRT opens the
possibility of \emph{continuous assay}: rather than testing at discrete
release points, the system continuously evaluates incoming production
traces against the behavioral baseline, raising a statistical alarm when
the fingerprint distribution shifts beyond a threshold. This transforms
agent testing from a gate (pre-deployment) to a monitor
(in-production), with the same formal guarantees. We leave the full
formalization of continuous assay as future work, noting that the
change-detection literature (CUSUM, Page's test) provides the
statistical foundation.

\subsection{Future Work}
\label{sec:discussion:future}

\paragraph{Compositional Testing.}
Our current framework tests individual agents and simple pipelines.
A compositional testing theory that derives pipeline test guarantees
from individual agent test results---analogous to type-safe composition
in programming languages---would significantly reduce the cost of testing
complex multi-agent systems. The composition conditions (C1--C4) from
the ABC framework~\citep{bhardwaj2026abc} provide a starting point.

\paragraph{Property-Based Test Generation.}
Integrating property-based testing~\citep{claessen2000quickcheck} with
stochastic test semantics would enable automatic generation of test
inputs from formal specifications. Given a contract $\phi$ and a
generator for inputs in $\mathcal{X}$, the framework could automatically
synthesize scenarios that stress-test the agent's contract compliance.

\paragraph{Online Regression Detection.}
Our framework performs offline testing (pre-deployment). Extending it to
online (production) regression detection---continuously monitoring agent
quality and triggering alerts when behavior degrades---would complement
the runtime enforcement capabilities of
\agentassert{}~\citep{bhardwaj2026abc}.

\paragraph{Multi-Agent System Testing.}
Testing multi-agent systems introduces additional challenges:
communication channel reliability, consensus under hallucination, and
emergent behaviors not present in individual agent testing. Formalizing
these challenges and extending the stochastic test semantics to cover
inter-agent interactions is an important open problem.

\paragraph{Human-in-the-Loop Testing.}
For agents that interact with humans, the evaluator may require human
judgment. Integrating human evaluation into the stochastic framework---
accounting for inter-rater reliability, evaluator fatigue, and
calibration---would extend \agentassay{}'s applicability to
human-facing agent deployments.
